% Fragment for the degree-ten C_2 x C_36 exclusion. No irreducibility is assumed.
\begin{lemma}\label{lem:236-degree10-oldN8}
Let $X=X_1(2,36)$, and let $\overline N$ be the noncuspidal
closed point of degree $8$ over $\mathbb F_5$ specified below.
There is no effective divisor $D$ on $X$, defined over $\mathbb Q$,
of degree $10$, whose support is entirely noncuspidal and whose
reduction at $5$ is $\overline N+\overline C$, where
$\overline C$ is a cuspidal effective divisor of degree $2$.
\end{lemma}
\begin{proof}
Put $g(P,Q)=(P+18Q,P+Q)$.
This automorphism has order $4$ and acts freely on the geometric
cusps. The point $\overline N$ is invariant under $g$.
Write $\pi:X\longrightarrow W=X/\langle g\rangle$.
The curve $W$ has genus $11$, the map $\pi$ is unramified, and
$\overline N=\pi^*q$, where $q$ has degree $2$ over $\mathbb F_5$.

We first compute
\begin{equation}\label{eq:236-degree10-cusp-section}
 h^0\bigl(\overline N+\pi^*\pi_*\overline C\bigr)=1
\end{equation}
for every cuspidal effective divisor $\overline C$ of degree $2$.
Here is a model of $W$ and the required calculation.
Put $r=t^2-t+1$, $c=t(r-1)$, $b=rc$, and
\[
 E_9:\quad y^2+(1-c)xy-by=x^3-bx^2.
\]
Then $W$ is given by
\[
 16(w^2+16)^3=j(E_9)w^4.
\]
After completing the square in $E_9$, write its equation as
$y^2=x^3+a_2x^2+a_4x+a_6$.
The cubic has a unique common root $\rho$ with
\[
 w^2(a_2+3x)^2-4(w^2+4)(a_4+2a_2x+3x^2).
\]
Put $\beta=a_4+2a_2\rho+3\rho^2$ and
$f=16\beta^3(1+4/w)^2$.
Over $\mathbb F_5$ the cover is $z^4=f$.
Its divisor is $\operatorname{div}(f)=4D_0$, where $[D_0]$ has
order $4$. Hence, for every divisor $A$ on $W$,
\[
 h^0_X(\pi^*A)=\sum_{j=0}^3h^0_W(A+jD_0).
\]

There are six rational cusps $c_1,\ldots,c_6$ of $W$ whose
fibres split into four rational cusps of $X$.
There are two quadratic cusps $u_1,u_2$ whose fibres split into
four quadratic cusps. These account for all source cusps of
degree at most $2$.
The exact Riemann--Roch calculations give
\[
 \bigl(h^0(A+jD_0)\bigr)_{j=0}^3=(1,0,0,0)
\]
for the two divisors $A=q+u_i$ and the $21$ divisors
$A=q+c_i+c_j$, where $1\leq i\leq j\leq6$.
This proves \eqref{eq:236-degree10-cusp-section}.

Suppose that $D$ exists. Work over a finite unramified extension
of $\mathbb Q_5$ containing the cusps, and let $C$ be the cuspidal
lift of $\overline C$. Put $L=[D-C]\in\operatorname{Pic}^8(X)$.
The Jacobian $J_X(\mathbb Q)$ has rank zero. Together with
Manin--Drinfeld, this implies that $gL-L$ is torsion.
Its reduction is zero, since $\overline N$ is $g$-invariant.
Torsion reduction is injective over an unramified extension of
$\mathbb Q_5$, and hence $gL=L$.

Put $\Omega=\pi^*\pi_*C=\sum_{j=0}^3g^jC$.
The divisors
\[
 D+\Omega-C\quad\hbox{and}\quad gD+\Omega-gC
\]
are effective and linearly equivalent.
Their common line bundle reduces to
$\mathcal O(\overline N+\pi^*\pi_*\overline C)$.
By \eqref{eq:236-degree10-cusp-section} and upper semicontinuity,
its space of sections has dimension at most $1$.
Thus the two divisors are equal.
Since $D$ and $gD$ have no cuspidal components, comparison of
cuspidal parts gives $C=gC$.
Every geometric cusp orbit has length $4$, whereas $\deg C=2$.
This is a contradiction.
\end{proof}

\noindent\emph{Identification and computational evidence.}
For an exact specification of $\overline N$, let
\[
 k=\mathbb F_5[\alpha]/
 (\alpha^8+\alpha^4+3\alpha^2+4\alpha+2).
\]
Write $B=\sum_{i=0}^7B_i\alpha^i$ and
$C=\sum_{i=0}^7C_i\alpha^i$, with coefficient vectors
\[
 (B_i)=(0,0,0,0,2,3,0,4),\qquad
 (C_i)=(2,1,0,4,0,2,1,2).
\]
The level structure
$E:y^2=x^3+Cx^2+B(1-B-C)x$, $P=(0,0)$, $Q=(B,B)$
represents $\overline N$. On it, $g$ agrees with the sixth
power of the $5$-power Frobenius.
The function-field calculation and all $23$ Riemann--Roch tests
are in \texttt{new236\_geom\_cusp2\_numeric.m}; the exported data
are in \texttt{geometry/cusp2\_numeric\_certificate.m}.
An independent exact replay in Sage verifies the level structure,
the Frobenius relation, the source-to-quotient invariants, the
local Kummer splitting, and the free action on all $64$ geometric
cusps; see \texttt{geometry/replay\_cusp2\_v2.py} and
\texttt{geometry/cusp2\_independent\_replay.json}.
